\documentclass[10pt]{article}
\usepackage[T1]{fontenc}
\usepackage{lmodern,microtype,amsmath,amssymb,array,tabularx,booktabs}
\usepackage[margin=0.85in]{geometry}
\usepackage{xurl,hyperref}
\hypersetup{colorlinks=true,urlcolor=blue,pdftitle={Appendix A review -- Version 17}}
\setlength{\parindent}{0pt}
\setlength{\parskip}{5pt}
\newcolumntype{Y}{>{\raggedright\arraybackslash}X}
\newcommand{\Two}{\mathbf2}
\begin{document}
{\LARGE Appendix A: revision and review}\par
{\large Version 17, 31 October 2026}\par

Version 17 computes the Boolean operations on coproduct summands, expands
the verification of their identities, and defines conditionals by an explicit copairing. Every downloadable artifact has a new versioned filename. The standalone appendix is
Appendix A; its mathematical body is Appendix B of the assembled manuscript.
Both PDFs identify version 17 in their metadata and displayed text.

\section*{Boolean operations and conditionals}
Let $j_{\mathbf{t}\mathbf{f}}:1\to1+1+1+1$ be the third coproduct
injection. The definition of $\gamma$ gives
$\gamma j_{\mathbf{t}\mathbf{f}}=\langle\mathbf{t},\mathbf{f}\rangle$.
Consequently,
\[
\wedge_2\langle\mathbf{t},\mathbf{f}\rangle
=[\mathbf{f},\mathbf{f},\mathbf{f},\mathbf{t}]j_{\mathbf{t}\mathbf{f}}
=\mathbf{f},\qquad
\vee_2\langle\mathbf{t},\mathbf{f}\rangle
=[\mathbf{f},\mathbf{t},\mathbf{t},\mathbf{t}]j_{\mathbf{t}\mathbf{f}}
=\mathbf{t}.
\]
These are defining equations of the copairings. The appendix gives all four
component calculations. It then verifies the Boolean identities on the jointly
epimorphic coproduct injections of the relevant finite products, including
explicit cases for associativity and both distributive identities. The argument
uses labelled summands and does not require the two truth injections to be distinct.

Let $p:X\to\Two$. Write $[p=\mathbf{t}]\hookrightarrow X$ and
$[p=\mathbf{f}]\hookrightarrow X$ for its pullbacks along the respective
truth injections. Denote these inclusions by $i_{\mathbf{t}}$ and
$i_{\mathbf{f}}$. The coproduct hypotheses give the isomorphism
\[
\delta_p=[i_{\mathbf{t}},i_{\mathbf{f}}]:
[p=\mathbf{t}]+[p=\mathbf{f}]\xrightarrow{\cong}X.
\]
For branches $f,g:X\to Y$, the conditional is
\[
\mathsf{if}(p;f,g)=[f i_{\mathbf{t}},g i_{\mathbf{f}}]\delta_p^{-1}:X\to Y.
\]
It is the unique morphism $h:X\to Y$ satisfying
$h i_{\mathbf{t}}=f i_{\mathbf{t}}$ and
$h i_{\mathbf{f}}=g i_{\mathbf{f}}$.

The bold notation is also used in the later pullbacks defining the
capture-test subobject and the two absence-of-variable subobjects.
All conditional expressions retain their original branch order:
the first branch corresponds to $\mathbf{t}$ and the second to $\mathbf{f}$.

\section*{Review scope and findings}
A separate agent checked the copairing's source and target, the inverse
of the coproduct decomposition, both restriction equations, the branch order
in existing conditionals, and the later pullback notation. The expanded Boolean
calculation was checked separately. The finite-product conventions introduced
in version 16 are retained.

The four-agent review of version 14 remains a separate record: hostile
mathematical review, mathematical-writing review, six-section outline, and
reconciliation. These are not fresh full reviews of version 17. Their two
principal corrections---explicit copairings and the initial-algebra composite
in the substitution proof---are retained.

\newpage
\section*{Six-section outline}
\renewcommand{\arraystretch}{1.15}
\begin{tabularx}{\textwidth}{@{}p{0.08\textwidth}>{\raggedright\arraybackslash}p{0.32\textwidth}Y@{}}
\toprule
Section & Title & Function and dependencies\\\midrule
A.1 & Categorical hypotheses and notation & Ambient structure, Boolean
operations and their laws, variable equality, signatures, and parameter stability.\\
A.2 & Terms, formulas, and free variables & Initial syntax and list algebras,
free-variable morphisms, and the sentence subobject; uses A.1.\\
A.3 & Admissible substitution & Single-variable substitution, its admissibility
subobject, and substitution laws; uses A.1--A.2.\\
A.4 & Inference rules and deductively closed subobjects & Inference endofunctor
and the least closed image theorem; needs regularity and finite coproducts.\\
A.5 & First-order theories & Logical schemes, formula and sentence theories,
and partial identities; uses A.2--A.4.\\
A.6 & Propositional theories & The propositional instance of A.4 under weaker
hypotheses, reusing the three propositional schemes of A.5.\\\bottomrule
\end{tabularx}

The section order, theorem labels, bibliography, and dependencies are unchanged.
Numerical realization proofs remain in the applications. The appendix constructs
the theory generated by an axiom subobject under its stated assumptions.

\section*{Publication and verification}
The deliverables have distinct names ending in \texttt{\_v17}: the standalone
PDF and LaTeX source, this review PDF and its source, the assembled manuscript
PDF, and its source archive. The current links point to these versioned files.
The standalone PDF also displays its version immediately below the title.
Earlier repository versions remain available.

The final sources were compiled, checked for unresolved references and layout
warnings, and visually inspected on the revised pages. The standalone and
integrated mathematical bodies agree after excluding their publication metadata
and the integrated compatibility label. The accepted abstract and other
mathematical sections are unchanged.

The version 17 check is limited to these constructions and their notation. No new full-manuscript, novelty,
bibliography, or Lean verification is claimed. The complete version 14 reports
remain at
\href{https://github.com/flengyel/CategoricalEssentialIncompleteness/tree/main/Audit/2026-10-31-appendix-v14}{the version 14 audit directory}.
\end{document}
